%==============================================================================
%  Linear Regression: Theory, Mathematical Foundations,
%  Closed-Form Solution and Gradient Descent
%
%  Self-contained lecture notes for an introductory Machine Learning course
%  (B.Tech / M.Tech level).  Prerequisites: basic calculus and linear algebra.
%
%  Compile with:  pdflatex linear_regression_notes.tex   (run twice for refs)
%  Requires a full TeX distribution (TeX Live / MiKTeX) with pgfplots.
%==============================================================================
\documentclass[11pt,a4paper]{report}

%------------------------------- Page geometry --------------------------------
\usepackage[a4paper,margin=1in]{geometry}

%------------------------------- Mathematics ----------------------------------
\usepackage{amsmath,amssymb,amsthm,mathtools}
\usepackage{bm}                 % bold math symbols

%------------------------------- Tables & layout ------------------------------
\usepackage{booktabs}
\usepackage{array}
\usepackage{multirow}
\usepackage{longtable}
\usepackage{enumitem}
\usepackage{caption}

%------------------------------- Graphics -------------------------------------
\usepackage{graphicx}
\usepackage{xcolor}
\usepackage{tikz}
\usepackage{pgfplots}
\pgfplotsset{compat=1.18}
\usetikzlibrary{arrows.meta,calc}

%------------------------------- Code listings --------------------------------
\usepackage{listings}

%------------------------------- Hyperlinks (load last) -----------------------
\usepackage[colorlinks=true,linkcolor=blue!60!black,citecolor=blue!60!black,
urlcolor=blue!60!black,linktocpage=true]{hyperref}

%------------------------------- Theorem-like envs ----------------------------
\theoremstyle{definition}
\newtheorem{definition}{Definition}[chapter]
\newtheorem{example}{Example}[chapter]
\newtheorem{exercise}{Exercise}[chapter]
\theoremstyle{plain}
\newtheorem{theorem}{Theorem}[chapter]
\newtheorem{proposition}{Proposition}[chapter]
\theoremstyle{remark}
\newtheorem*{remark}{Remark}

%------------------------------- Handy macros ---------------------------------
\newcommand{\R}{\mathbb{R}}
\newcommand{\E}{\mathbb{E}}
\newcommand{\vx}{\mathbf{x}}
\newcommand{\vy}{\mathbf{y}}
\newcommand{\vw}{\bm{\theta}}          % parameter vector theta
\newcommand{\Xmat}{\mathbf{X}}
\newcommand{\xbar}{\bar{x}}
\newcommand{\ybar}{\bar{y}}
\newcommand{\htheta}{h_{\bm{\theta}}}
\newcommand{\normeq}{\Xmat^{\!\top}\Xmat\,\vw = \Xmat^{\!\top}\vy}
\DeclareMathOperator{\Cov}{Cov}
\DeclareMathOperator{\Var}{Var}
\DeclareMathOperator*{\argmin}{arg\,min}

%------------------------------- Listing styles -------------------------------
\definecolor{codegreen}{rgb}{0.0,0.5,0.0}
\definecolor{codegray}{rgb}{0.5,0.5,0.5}
\definecolor{codeblue}{rgb}{0.1,0.1,0.6}
\definecolor{backcol}{rgb}{0.97,0.97,0.95}

\lstdefinestyle{common}{
	backgroundcolor=\color{backcol},
	commentstyle=\color{codegreen},
	keywordstyle=\color{codeblue}\bfseries,
	numberstyle=\tiny\color{codegray},
	stringstyle=\color{codegreen},
	basicstyle=\ttfamily\footnotesize,
	breakatwhitespace=false,
	breaklines=true,
	captionpos=b,
	keepspaces=true,
	numbers=left,
	numbersep=6pt,
	showspaces=false,
	showstringspaces=false,
	showtabs=false,
	frame=single,
	rulecolor=\color{codegray},
	tabsize=2
}
\lstdefinestyle{matlab}{style=common,language=Matlab}
\lstdefinestyle{python}{style=common,language=Python}

%==============================================================================
\begin{document}
	
	%------------------------------- Title ----------------------------------------
	\begin{titlepage}
		\centering
		\vspace*{2cm}
		{\Huge\bfseries Linear Regression\par}
		\vspace{0.6cm}
		{\Large Theory, Mathematical Foundations,\\[2pt]
			Closed-Form Solution and Gradient Descent\par}
		\vspace{1.5cm}
		{\large A Self-Contained Set of Lecture Notes\par}
		\vspace{0.4cm}
		{\large Introduction to Machine Learning\par}
		\vspace{2.5cm}
		\begin{tabular}{c}
			\hline\\[-6pt]
			\textbf{Prerequisites:} Basic Calculus and Linear Algebra\\[4pt]
			\textbf{Level:} B.Tech / M.Tech\\[4pt]
			\hline
		\end{tabular}
		\vfill
		{\large \today\par}
	\end{titlepage}
	
	\tableofcontents
	\clearpage
	
	%==============================================================================
	\chapter{Introduction to Machine Learning}
	%==============================================================================
	
	\section{What is Artificial Intelligence?}
	\emph{Artificial Intelligence} (AI) is the branch of computer science that
	studies how to build systems capable of performing tasks that normally require
	human intelligence---perception, reasoning, planning, language understanding
	and decision making. Classical AI encoded knowledge as explicit rules written
	by human experts. This works well in narrow, well-understood domains, but
	becomes brittle when the number of rules and exceptions explodes.
	
	\section{What is Machine Learning?}
	\emph{Machine Learning} (ML) is the sub-field of AI in which a system
	\emph{learns} the rules directly from data rather than having them
	hand-coded. Formally, following Tom Mitchell:
	
	\begin{quote}
		A computer program is said to learn from experience $E$ with respect to some
		class of tasks $T$ and performance measure $P$, if its performance at tasks in
		$T$, as measured by $P$, improves with experience $E$.
	\end{quote}
	
	For linear regression the ingredients are:
	\begin{itemize}[nosep]
		\item $T$: predict a real-valued output from input features;
		\item $E$: a dataset of past $(\text{input},\text{output})$ pairs;
		\item $P$: a loss (error) function such as the mean squared error.
	\end{itemize}
	
	\section{Types of Machine Learning}
	
	\subsection{Supervised Learning}
	The training data consists of input--output pairs
	$\{(\vx^{(i)},y^{(i)})\}_{i=1}^{m}$, and the goal is to learn a function
	$h:\mathcal{X}\to\mathcal{Y}$ that predicts $y$ from $\vx$. When $y$ is
	continuous the problem is called \emph{regression}; when $y$ is a discrete
	label it is called \emph{classification}. \textbf{Linear regression is the
		canonical supervised regression algorithm.}
	
	\subsection{Unsupervised Learning}
	Only inputs $\{\vx^{(i)}\}$ are available (no labels). The goal is to discover
	structure---clusters, low-dimensional manifolds, density. Examples: $k$-means,
	principal component analysis.
	
	\subsection{Reinforcement Learning}
	An agent interacts with an environment, receives rewards, and learns a policy
	that maximizes cumulative reward over time. There is no fixed labelled dataset;
	feedback is evaluative and delayed.
	
	\section{Where does Linear Regression fit?}
	Linear regression is \emph{supervised}, solves a \emph{regression} task, and
	uses a \emph{parametric linear} model. Despite its simplicity it is the
	foundation for a large part of statistics and ML: logistic regression,
	generalized linear models, ridge/lasso regularization, and even the final
	layer of many neural networks are direct descendants.
	
	\section{Applications of Linear Regression}
	\begin{itemize}[nosep]
		\item Predicting house prices from size, location and number of rooms.
		\item Forecasting sales or demand from advertising spend.
		\item Estimating a person's expenditure from income.
		\item Modelling the relationship between temperature and energy consumption.
		\item Calibrating sensors (mapping raw readings to physical quantities).
	\end{itemize}
	
	\section{Regression vs Classification}
	\begin{center}
		\begin{tabular}{@{}p{6.5cm}p{6.5cm}@{}}
			\toprule
			\textbf{Regression} & \textbf{Classification}\\
			\midrule
			Output is continuous ($y\in\R$).      & Output is discrete (a class label).\\
			Example: price, temperature.          & Example: spam / not-spam.\\
			Typical loss: squared error.          & Typical loss: cross-entropy.\\
			Metric: RMSE, $R^2$.                  & Metric: accuracy, F1.\\
			\bottomrule
		\end{tabular}
	\end{center}
	
	\section{Assumptions of Linear Regression}
	The ordinary least-squares (OLS) model relies on the following assumptions.
	\begin{enumerate}[nosep]
		\item \textbf{Linearity.} The expected value of $y$ is a linear function of
		the parameters.
		\item \textbf{Independence.} The observations (errors) are independent.
		\item \textbf{Homoscedasticity.} The errors have constant variance
		$\sigma^2$.
		\item \textbf{Normality.} For inference, errors are (approximately) normally
		distributed with zero mean.
		\item \textbf{No perfect multicollinearity.} The feature matrix has full
		column rank, so $\Xmat^{\!\top}\Xmat$ is invertible.
	\end{enumerate}
	
	%==============================================================================
	\chapter{Mathematical Background}
	%==============================================================================
	
	\section{Scalars, Vectors and Matrices}
	A \emph{scalar} is a single number, e.g.\ $a\in\R$. A \emph{vector} is an
	ordered list of numbers, written as a column
	\[
	\vx=\begin{bmatrix}x_1\\x_2\\\vdots\\x_n\end{bmatrix}\in\R^{n}.
	\]
	A \emph{matrix} is a rectangular array $\Xmat\in\R^{m\times n}$ with entries
	$X_{ij}$. Transposition swaps rows and columns: $(\Xmat^{\!\top})_{ij}=X_{ji}$.
	
	\section{Dataset Representation}
	A supervised dataset with $m$ examples and $n$ features is written
	$\{(\vx^{(i)},y^{(i)})\}_{i=1}^{m}$, where $\vx^{(i)}\in\R^{n}$ and
	$y^{(i)}\in\R$. The superscript $(i)$ indexes examples, the subscript $j$
	indexes features, so $x^{(i)}_j$ is the $j$-th feature of the $i$-th example.
	
	\subsection{Feature Matrix}
	Stacking all examples as rows and prepending a column of ones (for the
	intercept term) gives the \emph{design} or \emph{feature matrix}
	\begin{equation}
		\Xmat=
		\begin{bmatrix}
			1 & x^{(1)}_1 & \cdots & x^{(1)}_n\\
			1 & x^{(2)}_1 & \cdots & x^{(2)}_n\\
			\vdots & \vdots & & \vdots\\
			1 & x^{(m)}_1 & \cdots & x^{(m)}_n
		\end{bmatrix}\in\R^{m\times(n+1)}.
		\label{eq:designmatrix}
	\end{equation}
	
	\subsection{Target Vector}
	The outputs are collected into the \emph{target vector}
	\[
	\vy=\begin{bmatrix}y^{(1)}\\y^{(2)}\\\vdots\\y^{(m)}\end{bmatrix}\in\R^{m}.
	\]
	
	\section{Dot Product}
	For $\mathbf{a},\mathbf{b}\in\R^{n}$,
	\[
	\mathbf{a}^{\!\top}\mathbf{b}=\sum_{j=1}^{n}a_j b_j .
	\]
	Geometrically $\mathbf{a}^{\!\top}\mathbf{b}=\|\mathbf{a}\|\,\|\mathbf{b}\|\cos\vartheta$.
	
	\section{Matrix Multiplication}
	If $\mathbf{A}\in\R^{m\times k}$ and $\mathbf{B}\in\R^{k\times n}$ then
	$\mathbf{C}=\mathbf{A}\mathbf{B}\in\R^{m\times n}$ with
	$C_{ij}=\sum_{\ell=1}^{k}A_{i\ell}B_{\ell j}$. In particular, the prediction of
	a linear model for all examples at once is the matrix--vector product
	$\Xmat\vw$.
	
	\section{Descriptive Statistics}
	Let $x_1,\dots,x_m$ be a sample.
	
	\begin{definition}[Mean]
		\[
		\xbar=\frac1m\sum_{i=1}^{m}x_i .
		\]
	\end{definition}
	
	\begin{definition}[Variance]
		The (population) variance measures spread about the mean:
		\[
		\Var(x)=\frac1m\sum_{i=1}^{m}(x_i-\xbar)^2 .
		\]
		The unbiased sample variance divides by $m-1$ instead of $m$.
	\end{definition}
	
	\begin{definition}[Standard deviation]
		\[
		\sigma_x=\sqrt{\Var(x)} .
		\]
	\end{definition}
	
	\begin{definition}[Covariance]
		For paired data $(x_i,y_i)$,
		\[
		\Cov(x,y)=\frac1m\sum_{i=1}^{m}(x_i-\xbar)(y_i-\ybar).
		\]
		Covariance is positive when the two variables tend to increase together.
	\end{definition}
	
	\begin{definition}[Correlation coefficient]
		The Pearson correlation coefficient normalizes covariance to $[-1,1]$:
		\begin{equation}
			r=\frac{\Cov(x,y)}{\sigma_x\,\sigma_y}
			=\frac{\sum_{i}(x_i-\xbar)(y_i-\ybar)}
			{\sqrt{\sum_{i}(x_i-\xbar)^2}\,\sqrt{\sum_{i}(y_i-\ybar)^2}}.
			\label{eq:corr}
		\end{equation}
	\end{definition}
	
	These five quantities reappear directly in the closed-form solution of simple
	linear regression (Chapter~\ref{ch:closedform}).
	
	%==============================================================================
	\chapter{Linear Regression Theory}
	%==============================================================================
	
	\section{Problem Formulation}
	Given training data $\{(\vx^{(i)},y^{(i)})\}_{i=1}^{m}$ we seek a function that
	predicts $y$ from $\vx$ as accurately as possible on unseen data.
	
	\paragraph{Training data.} The $m$ observed pairs used to fit the model.
	\paragraph{Features.} The inputs $\vx^{(i)}\in\R^{n}$ (also called predictors,
	covariates or independent variables).
	\paragraph{Target.} The output $y^{(i)}\in\R$ (also called the response or
	dependent variable).
	
	\section{The Linear Model}
	Linear regression assumes the target is (approximately) a linear function of
	the features plus an intercept:
	\begin{equation}
		\htheta(\vx)=\theta_0+\theta_1 x_1+\theta_2 x_2+\cdots+\theta_n x_n .
		\label{eq:multimodel}
	\end{equation}
	Introducing a dummy feature $x_0\equiv 1$ this becomes the compact dot product
	\begin{equation}
		\htheta(\vx)=\sum_{j=0}^{n}\theta_j x_j=\vw^{\!\top}\vx ,
		\qquad \vw=\begin{bmatrix}\theta_0\\\vdots\\\theta_n\end{bmatrix},\;
		\vx=\begin{bmatrix}1\\x_1\\\vdots\\x_n\end{bmatrix}.
		\label{eq:dotmodel}
	\end{equation}
	
	\subsection{Simple Linear Regression}
	The special case $n=1$ (a single feature) is \emph{simple linear regression}:
	\begin{equation}
		\boxed{\;\hat{y}=\htheta(x)=\theta_0+\theta_1 x\;}
		\label{eq:simplemodel}
	\end{equation}
	Here $\theta_0$ is the \emph{intercept} (value of $\hat y$ when $x=0$) and
	$\theta_1$ is the \emph{slope} (change in $\hat y$ per unit change in $x$).
	
	\paragraph{Hypothesis function and prediction.}
	The function $\htheta$ is the \emph{hypothesis}; evaluating it at a new input
	$x_\star$ yields the \emph{prediction} $\hat y_\star=\theta_0+\theta_1 x_\star$.
	Learning means choosing $\vw$ so that predictions match the targets.
	
	%==============================================================================
	\chapter{The Loss Function}
	%==============================================================================
	
	\section{Definition}
	To choose $\vw$ we need a measure of how wrong the predictions are. Define the
	error (residual) of example $i$ as $\htheta(x^{(i)})-y^{(i)}$. The \emph{mean
		squared error} cost function is
	\begin{equation}
		\boxed{\;J(\vw)=\frac{1}{2m}\sum_{i=1}^{m}\bigl(\htheta(x^{(i)})-y^{(i)}\bigr)^2\;}
		\label{eq:cost}
	\end{equation}
	
	\section{Why Square the Error?}
	\begin{itemize}[nosep]
		\item \textbf{Sign removal.} Squaring makes positive and negative errors both
		count as ``bad'' and prevents them from cancelling.
		\item \textbf{Smoothness.} $(\cdot)^2$ is differentiable everywhere, unlike
		$|\cdot|$, which lets us use calculus and gradient descent.
		\item \textbf{Penalizing large errors.} Big mistakes are penalized
		disproportionately, encouraging the fit to avoid gross outliers.
		\item \textbf{Statistical justification.} Under Gaussian noise, minimizing
		squared error is equivalent to maximum-likelihood estimation.
	\end{itemize}
	
	\section{Why Divide by $m$ (and by $2$)?}
	Dividing by $m$ turns the \emph{sum} of squared errors into a \emph{mean},
	making the cost independent of the dataset size so that learning rates transfer
	across datasets. The extra factor of $\tfrac12$ is a convenience: when we
	differentiate the square, the exponent $2$ cancels the $\tfrac12$, giving a
	cleaner gradient. Neither constant changes the location of the minimum:
	\[
	\argmin_{\vw} \tfrac{1}{2m}\textstyle\sum_i(\cdot)^2
	= \argmin_{\vw} \sum_i(\cdot)^2 .
	\]
	
	\section{Convexity of the Cost Function}
	\label{sec:convex}
	As a function of $\vw$, $J$ is a sum of squares of \emph{affine} functions of
	$\vw$, hence a quadratic form. Its Hessian (Chapter~\ref{ch:normaleq}) is
	\[
	\nabla^2 J(\vw)=\frac1m\Xmat^{\!\top}\Xmat ,
	\]
	which is positive semi-definite because for any vector $\mathbf{v}$,
	$\mathbf{v}^{\!\top}\Xmat^{\!\top}\Xmat\mathbf{v}=\|\Xmat\mathbf{v}\|^2\ge0$.
	Therefore $J$ is \textbf{convex}: any stationary point is a \emph{global}
	minimum, and there are no spurious local minima. If $\Xmat$ has full column
	rank the Hessian is positive definite and the minimizer is unique.
	
	\begin{figure}[ht]
		\centering
		\begin{tikzpicture}
			\begin{axis}[
				width=10cm,height=7cm,
				view={35}{30},
				xlabel={$\theta_0$},ylabel={$\theta_1$},zlabel={$J(\bm\theta)$},
				colormap/viridis, grid=both,
				domain=-2:2,y domain=-2:2,samples=25,
				title={Convex ``bowl'' cost surface}]
				\addplot3[surf,opacity=0.85] {0.5*(x*x + y*y) + 0.3*x*y};
			\end{axis}
		\end{tikzpicture}
		\caption{The mean-squared-error cost is a convex paraboloid in parameter
			space; gradient descent slides down to its single global minimum.}
		\label{fig:bowl}
	\end{figure}
	
	%==============================================================================
	\chapter{Closed-Form Solution (Simple Linear Regression)}
	\label{ch:closedform}
	%==============================================================================
	
	We derive the optimal $\theta_0,\theta_1$ for the single-feature model
	\eqref{eq:simplemodel} by minimizing the sum of squared errors
	\begin{equation}
		S(\theta_0,\theta_1)=\sum_{i=1}^{m}\bigl(y^{(i)}-\theta_0-\theta_1 x^{(i)}\bigr)^2 .
		\label{eq:SSE}
	\end{equation}
	(We minimize $S=2m\,J$; the constant does not affect the minimizer.) To keep
	the notation light, write $x_i,y_i$ for $x^{(i)},y^{(i)}$ and let all sums run
	$i=1,\dots,m$.
	
	\section{Setting the Partial Derivatives to Zero}
	At a minimum both partial derivatives vanish.
	
	\paragraph{With respect to $\theta_0$.}
	\begin{align}
		\frac{\partial S}{\partial\theta_0}
		&=\sum_i 2\bigl(y_i-\theta_0-\theta_1 x_i\bigr)(-1)
		=-2\sum_i\bigl(y_i-\theta_0-\theta_1 x_i\bigr).
	\end{align}
	Setting $\dfrac{\partial S}{\partial\theta_0}=0$ and dividing by $-2$:
	\begin{equation}
		\sum_i\bigl(y_i-\theta_0-\theta_1 x_i\bigr)=0
		\;\Longrightarrow\;
		\sum_i y_i = m\,\theta_0+\theta_1\sum_i x_i .
		\label{eq:normal1}
	\end{equation}
	
	\paragraph{With respect to $\theta_1$.}
	\begin{align}
		\frac{\partial S}{\partial\theta_1}
		&=\sum_i 2\bigl(y_i-\theta_0-\theta_1 x_i\bigr)(-x_i)
		=-2\sum_i x_i\bigl(y_i-\theta_0-\theta_1 x_i\bigr).
	\end{align}
	Setting $\dfrac{\partial S}{\partial\theta_1}=0$ and dividing by $-2$:
	\begin{equation}
		\sum_i x_i y_i = \theta_0\sum_i x_i+\theta_1\sum_i x_i^2 .
		\label{eq:normal2}
	\end{equation}
	
	Equations \eqref{eq:normal1} and \eqref{eq:normal2} are the \textbf{two normal
		equations} for simple linear regression.
	
	\section{Solving the Normal Equations}
	From \eqref{eq:normal1}, dividing by $m$ and using
	$\xbar=\frac1m\sum x_i,\ \ybar=\frac1m\sum y_i$:
	\begin{equation}
		\ybar=\theta_0+\theta_1\xbar
		\;\Longrightarrow\;
		\boxed{\;\theta_0=\ybar-\theta_1\xbar\;}
		\label{eq:theta0}
	\end{equation}
	This shows the fitted line always passes through the centroid
	$(\xbar,\ybar)$. Substitute \eqref{eq:theta0} into \eqref{eq:normal2}:
	\begin{align}
		\sum_i x_i y_i
		&=(\ybar-\theta_1\xbar)\sum_i x_i+\theta_1\sum_i x_i^2\\
		&=\ybar\sum_i x_i+\theta_1\Bigl(\sum_i x_i^2-\xbar\sum_i x_i\Bigr).
	\end{align}
	Hence
	\begin{equation}
		\theta_1=\frac{\sum_i x_i y_i-\ybar\sum_i x_i}
		{\sum_i x_i^2-\xbar\sum_i x_i}.
		\label{eq:theta1raw}
	\end{equation}
	Using the algebraic identities
	$\sum_i x_i y_i-\ybar\sum_i x_i=\sum_i(x_i-\xbar)(y_i-\ybar)$ and
	$\sum_i x_i^2-\xbar\sum_i x_i=\sum_i(x_i-\xbar)^2$ (proved below), we obtain the
	standard form
	\begin{equation}
		\boxed{\;\theta_1
			=\frac{\sum_i(x_i-\xbar)(y_i-\ybar)}{\sum_i(x_i-\xbar)^2}
			=\frac{S_{xy}}{S_{xx}}\;}
		\label{eq:theta1}
	\end{equation}
	where $S_{xy}=\sum_i(x_i-\xbar)(y_i-\ybar)$ and
	$S_{xx}=\sum_i(x_i-\xbar)^2$.
	
	\begin{proof}[Proof of the two identities]
		Expand and use $\sum_i x_i=m\xbar,\ \sum_i y_i=m\ybar$:
		\begin{align*}
			\sum_i(x_i-\xbar)(y_i-\ybar)
			&=\sum_i x_i y_i-\ybar\sum_i x_i-\xbar\sum_i y_i+m\xbar\ybar\\
			&=\sum_i x_i y_i-\ybar(m\xbar)-\xbar(m\ybar)+m\xbar\ybar
			=\sum_i x_i y_i-m\xbar\ybar
			=\sum_i x_i y_i-\ybar\sum_i x_i .
		\end{align*}
		The second identity is the same computation with $y\equiv x$. \qedhere
	\end{proof}
	
	\section{Equivalent Forms via Covariance, Variance and Correlation}
	Dividing numerator and denominator of \eqref{eq:theta1} by $m$ converts the
	sums-of-squares into the statistics of Chapter~2:
	\begin{equation}
		\theta_1=\frac{\frac1m S_{xy}}{\frac1m S_{xx}}
		=\frac{\Cov(x,y)}{\Var(x)} .
		\label{eq:theta1cov}
	\end{equation}
	Furthermore, using $r=\dfrac{\Cov(x,y)}{\sigma_x\sigma_y}$ and
	$\Var(x)=\sigma_x^2$,
	\begin{equation}
		\boxed{\;\theta_1=\frac{\Cov(x,y)}{\Var(x)}
			=\frac{r\,\sigma_x\sigma_y}{\sigma_x^2}
			=r\,\frac{\sigma_y}{\sigma_x}\;}
		\label{eq:theta1r}
	\end{equation}
	So the slope equals the correlation coefficient scaled by the ratio of standard
	deviations, and the intercept follows from \eqref{eq:theta0}. These are the
	formulas we will use for hand computation in Chapter~\ref{ch:cfexample}.
	
	%==============================================================================
	\chapter{Matrix Formulation}
	%==============================================================================
	
	We now treat the general multi-feature model in matrix notation, which yields a
	single elegant formula valid for any number of features.
	
	\section{Model in Matrix Form}
	With the design matrix $\Xmat\in\R^{m\times(n+1)}$ of \eqref{eq:designmatrix}
	(a leading column of ones) and parameter vector
	$\vw=[\theta_0,\theta_1,\dots,\theta_n]^{\!\top}$, the vector of all $m$
	predictions is
	\begin{equation}
		\hat{\vy}=\Xmat\vw .
		\label{eq:matpred}
	\end{equation}
	For the single-feature case,
	\[
	\Xmat=\begin{bmatrix}1&x_1\\1&x_2\\\vdots&\vdots\\1&x_m\end{bmatrix},\qquad
	\vw=\begin{bmatrix}\theta_0\\\theta_1\end{bmatrix},\qquad
	\hat{\vy}=\Xmat\vw=\begin{bmatrix}\theta_0+\theta_1 x_1\\\vdots\\\theta_0+\theta_1 x_m\end{bmatrix}.
	\]
	
	\section{Residual Vector and the Cost}
	The \emph{residual vector} collects all errors,
	\begin{equation}
		\mathbf{r}=\Xmat\vw-\vy
		=\begin{bmatrix}\htheta(x^{(1)})-y^{(1)}\\\vdots\\\htheta(x^{(m)})-y^{(m)}\end{bmatrix}.
		\label{eq:residual}
	\end{equation}
	The key observation is that the sum of squared errors is exactly the squared
	Euclidean norm of the residual vector, i.e.\ its dot product with itself:
	\begin{equation}
		\sum_{i=1}^{m}e_i^2
		=\sum_{i=1}^{m}\bigl(\htheta(x^{(i)})-y^{(i)}\bigr)^2
		=\mathbf{r}^{\!\top}\mathbf{r}
		=(\Xmat\vw-\vy)^{\!\top}(\Xmat\vw-\vy).
		\label{eq:sumtonorm}
	\end{equation}
	
	\begin{proof}[Why \eqref{eq:sumtonorm} holds --- no skipped steps]
		For \emph{any} vector $\mathbf{r}=[e_1,\dots,e_m]^{\!\top}\in\R^{m}$, the
		definition of matrix multiplication gives the dot product
		\[
		\mathbf{r}^{\!\top}\mathbf{r}
		=\begin{bmatrix}e_1&e_2&\cdots&e_m\end{bmatrix}
		\begin{bmatrix}e_1\\e_2\\\vdots\\e_m\end{bmatrix}
		=e_1e_1+e_2e_2+\cdots+e_me_m
		=\sum_{i=1}^{m}e_i^2 .
		\]
		Now substitute the specific residual $\mathbf{r}=\Xmat\vw-\vy$ from
		\eqref{eq:residual}; its $i$-th component is exactly
		$e_i=\htheta(x^{(i)})-y^{(i)}$, which proves \eqref{eq:sumtonorm}.
	\end{proof}
	
	Therefore the cost function \eqref{eq:cost} in matrix form is
	\begin{equation}
		\boxed{\;J(\vw)=\frac{1}{2m}(\Xmat\vw-\vy)^{\!\top}(\Xmat\vw-\vy)\;}
		\label{eq:matcost}
	\end{equation}
	
	%==============================================================================
	\chapter{Derivation of the Normal Equation}
	\label{ch:normaleq}
	%==============================================================================
	
	We minimize \eqref{eq:matcost} directly with matrix calculus. Because a
	constant scaling does not move the minimizer, drop the $\tfrac1m$ and work with
	$J(\vw)=\tfrac12(\Xmat\vw-\vy)^{\!\top}(\Xmat\vw-\vy)$.
	
	\section{Expanding the Quadratic Form}
	Let $\mathbf{r}=\Xmat\vw-\vy$. Then
	\begin{align}
		2J(\vw)&=(\Xmat\vw-\vy)^{\!\top}(\Xmat\vw-\vy)\notag\\
		&=(\Xmat\vw)^{\!\top}(\Xmat\vw)
		-(\Xmat\vw)^{\!\top}\vy
		-\vy^{\!\top}(\Xmat\vw)
		+\vy^{\!\top}\vy .
		\label{eq:expand1}
	\end{align}
	Using $(\Xmat\vw)^{\!\top}=\vw^{\!\top}\Xmat^{\!\top}$ and noting that each term
	is a $1\times1$ scalar (so it equals its own transpose), the two cross terms
	are identical:
	\[
	(\Xmat\vw)^{\!\top}\vy=\vw^{\!\top}\Xmat^{\!\top}\vy,\qquad
	\vy^{\!\top}(\Xmat\vw)=\bigl(\vw^{\!\top}\Xmat^{\!\top}\vy\bigr)^{\!\top}
	=\vw^{\!\top}\Xmat^{\!\top}\vy .
	\]
	Hence \eqref{eq:expand1} collapses to
	\begin{equation}
		2J(\vw)=\vw^{\!\top}\Xmat^{\!\top}\Xmat\vw
		-2\,\vw^{\!\top}\Xmat^{\!\top}\vy
		+\vy^{\!\top}\vy .
		\label{eq:expand2}
	\end{equation}
	
	\section{Matrix Calculus Rules Used}
	We use two standard vector-derivative identities. For constant vector
	$\mathbf{b}$ and constant symmetric matrix $\mathbf{A}=\mathbf{A}^{\!\top}$,
	\begin{equation}
		\frac{\partial}{\partial\vw}\bigl(\vw^{\!\top}\mathbf{b}\bigr)=\mathbf{b},
		\qquad
		\frac{\partial}{\partial\vw}\bigl(\vw^{\!\top}\mathbf{A}\vw\bigr)=2\mathbf{A}\vw .
		\label{eq:matcalc}
	\end{equation}
	
	\begin{proof}[Sketch of the second identity]
		Write $f(\vw)=\vw^{\!\top}\mathbf{A}\vw=\sum_{p}\sum_{q}\theta_p A_{pq}\theta_q$.
		Differentiate with respect to a single component $\theta_k$:
		\[
		\frac{\partial f}{\partial\theta_k}
		=\sum_q A_{kq}\theta_q+\sum_p \theta_p A_{pk}
		=(\mathbf{A}\vw)_k+(\mathbf{A}^{\!\top}\vw)_k .
		\]
		Stacking over $k$ gives $\nabla f=(\mathbf{A}+\mathbf{A}^{\!\top})\vw$, which
		equals $2\mathbf{A}\vw$ when $\mathbf{A}$ is symmetric. Here
		$\mathbf{A}=\Xmat^{\!\top}\Xmat$ is symmetric since
		$(\Xmat^{\!\top}\Xmat)^{\!\top}=\Xmat^{\!\top}\Xmat$.
	\end{proof}
	
	\section{Differentiating and Solving}
	Apply \eqref{eq:matcalc} term-by-term to \eqref{eq:expand2} (the constant
	$\vy^{\!\top}\vy$ has zero derivative):
	\begin{equation}
		\frac{\partial (2J)}{\partial\vw}
		=2\,\Xmat^{\!\top}\Xmat\vw-2\,\Xmat^{\!\top}\vy .
	\end{equation}
	Setting the gradient to zero,
	\begin{equation}
		\Xmat^{\!\top}\Xmat\vw-\Xmat^{\!\top}\vy=\mathbf{0}
		\;\Longrightarrow\;
		\boxed{\;\normeq\;}
		\label{eq:normaleq}
	\end{equation}
	which is the celebrated \textbf{normal equation}. If $\Xmat^{\!\top}\Xmat$ is
	invertible (full column rank, no perfect multicollinearity),
	\begin{equation}
		\boxed{\;\vw=(\Xmat^{\!\top}\Xmat)^{-1}\Xmat^{\!\top}\vy\;}
		\label{eq:normalsol}
	\end{equation}
	Because $J$ is convex (Section~\ref{sec:convex}), this stationary point is the
	unique global minimizer.
	
	\begin{remark}
		Equation \eqref{eq:normaleq} for the single-feature design matrix reproduces
		exactly the two scalar normal equations \eqref{eq:normal1}--\eqref{eq:normal2}:
		carrying out $\Xmat^{\!\top}\Xmat$ and $\Xmat^{\!\top}\vy$ gives
		\[
		\Xmat^{\!\top}\Xmat=\begin{bmatrix}m&\sum x_i\\[2pt]\sum x_i&\sum x_i^2\end{bmatrix},
		\qquad
		\Xmat^{\!\top}\vy=\begin{bmatrix}\sum y_i\\[2pt]\sum x_i y_i\end{bmatrix}.
		\]
	\end{remark}
	
	%==============================================================================
	\chapter{Gradient Descent}
	%==============================================================================
	
	\section{Why Iterative Optimization?}
	The closed form \eqref{eq:normalsol} requires inverting the
	$(n+1)\times(n+1)$ matrix $\Xmat^{\!\top}\Xmat$, which costs
	$\mathcal{O}(n^3)$ and needs the full matrix in memory. When the number of
	features $n$ (or examples $m$) is very large this becomes impractical.
	\emph{Gradient descent} avoids matrix inversion by iteratively stepping
	downhill on the cost surface (Figure~\ref{fig:bowl}).
	
	\section{The Gradient}
	The gradient of the cost \eqref{eq:cost} points in the direction of steepest
	increase; we move in the opposite direction. Differentiating $J$ with scalar
	calculus, exactly as in Chapter~\ref{ch:closedform} but keeping the $\tfrac1m$:
	\begin{align}
		\frac{\partial J}{\partial\theta_0}
		&=\frac1m\sum_{i=1}^{m}\bigl(\htheta(x^{(i)})-y^{(i)}\bigr),
		\label{eq:grad0}\\[4pt]
		\frac{\partial J}{\partial\theta_1}
		&=\frac1m\sum_{i=1}^{m}\bigl(\htheta(x^{(i)})-y^{(i)}\bigr)\,x^{(i)} .
		\label{eq:grad1}
	\end{align}
	For the general model the compact vector form is
	$\nabla J(\vw)=\frac1m\Xmat^{\!\top}(\Xmat\vw-\vy)$.
	
	\section{Update Equations}
	Starting from an initial guess, repeat until convergence (simultaneously
	updating all parameters):
	\begin{equation}
		\boxed{\;
			\begin{aligned}
				\theta_0 &:= \theta_0-\alpha\,\frac1m\sum_{i=1}^{m}\bigl(\htheta(x^{(i)})-y^{(i)}\bigr),\\[3pt]
				\theta_1 &:= \theta_1-\alpha\,\frac1m\sum_{i=1}^{m}\bigl(\htheta(x^{(i)})-y^{(i)}\bigr)\,x^{(i)},
			\end{aligned}\;}
		\label{eq:gdupdate}
	\end{equation}
	or, in vector form, $\vw:=\vw-\alpha\,\nabla J(\vw)$. The scalar $\alpha>0$ is
	the \emph{learning rate}.
	
	\section{Choice of Learning Rate}
	\begin{itemize}[nosep]
		\item If $\alpha$ is \textbf{too small}, convergence is very slow.
		\item If $\alpha$ is \textbf{too large}, the updates overshoot and $J$ may
		oscillate or diverge.
		\item A good practice is to plot $J$ versus iteration: it should decrease
		smoothly and monotonically.
		\item \textbf{Feature scaling} (standardizing each feature to zero mean and
		unit variance) makes the bowl of Figure~\ref{fig:bowl} more circular
		and dramatically speeds up convergence.
	\end{itemize}
	
	\section{Convex Optimization Guarantee}
	Because $J$ is convex, gradient descent with a sufficiently small constant
	$\alpha$ is guaranteed to converge to the \emph{global} minimum---the very same
	solution given by the normal equation \eqref{eq:normalsol}.
	
	%==============================================================================
	\chapter{Complete Worked Example: Closed Form}
	\label{ch:cfexample}
	%==============================================================================
	
	\section{The Dataset}
	Consider a tiny house-price dataset: $x$ is the size (in thousands of square
	feet) and $y$ is the price (in suitable monetary units). There are $m=5$
	examples.
	\begin{center}
		\begin{tabular}{@{}ccccccc@{}}
			\toprule
			$x$ (size) & 1 & 2 & 3 & 4 & 5 & $\sum=15$\\
			$y$ (price) & 1 & 3 & 2 & 5 & 4 & $\sum=15$\\
			\bottomrule
		\end{tabular}
	\end{center}
	
	\section{Step 1 --- Means}
	\[
	\xbar=\frac{1+2+3+4+5}{5}=\frac{15}{5}=3,\qquad
	\ybar=\frac{1+3+2+5+4}{5}=\frac{15}{5}=3 .
	\]
	
	\section{Step 2 --- Deviations, Covariance and Variance}
	\begin{center}
		\begin{tabular}{@{}ccccc@{}}
			\toprule
			$x_i$ & $y_i$ & $x_i-\xbar$ & $y_i-\ybar$ & $(x_i-\xbar)(y_i-\ybar)$ \\
			\midrule
			1 & 1 & $-2$ & $-2$ & $4$ \\
			2 & 3 & $-1$ & $ 0$ & $0$ \\
			3 & 2 & $ 0$ & $-1$ & $0$ \\
			4 & 5 & $ 1$ & $ 2$ & $2$ \\
			5 & 4 & $ 2$ & $ 1$ & $2$ \\
			\midrule
			\multicolumn{4}{r}{$S_{xy}=\sum(x_i-\xbar)(y_i-\ybar)=$} & $8$ \\
			\bottomrule
		\end{tabular}
	\end{center}
	Similarly
	\[
	S_{xx}=\sum(x_i-\xbar)^2=4+1+0+1+4=10,\qquad
	S_{yy}=\sum(y_i-\ybar)^2=4+0+1+4+1=10 .
	\]
	Hence
	\[
	\Cov(x,y)=\tfrac{S_{xy}}{m}=\tfrac{8}{5}=1.6,\qquad
	\Var(x)=\tfrac{S_{xx}}{m}=\tfrac{10}{5}=2,\qquad
	\sigma_x=\sigma_y=\sqrt{2}\approx1.414 .
	\]
	
	\section{Step 3 --- Slope and Intercept}
	Using \eqref{eq:theta1}/\eqref{eq:theta1cov},
	\[
	\theta_1=\frac{S_{xy}}{S_{xx}}=\frac{8}{10}=0.8
	=\frac{\Cov(x,y)}{\Var(x)}=\frac{1.6}{2}=0.8 .
	\]
	Cross-check with \eqref{eq:theta1r}: the correlation is
	$r=\dfrac{S_{xy}}{\sqrt{S_{xx}S_{yy}}}=\dfrac{8}{\sqrt{10\cdot10}}=0.8$, so
	$\theta_1=r\,\dfrac{\sigma_y}{\sigma_x}=0.8\cdot1=0.8$. \checkmark
	
	Then from \eqref{eq:theta0},
	\[
	\theta_0=\ybar-\theta_1\xbar=3-0.8\times3=3-2.4=0.6 .
	\]
	
	\section{Step 4 --- The Fitted Model}
	\begin{equation}
		\boxed{\;\hat y = 0.6 + 0.8\,x\;}
	\end{equation}
	Predictions, residuals and the sum of squared errors:
	\begin{center}
		\begin{tabular}{@{}cccc@{}}
			\toprule
			$x_i$ & $\hat y_i=0.6+0.8x_i$ & $e_i=y_i-\hat y_i$ & $e_i^2$\\
			\midrule
			1 & $1.4$ & $-0.4$ & $0.16$\\
			2 & $2.2$ & $ 0.8$ & $0.64$\\
			3 & $3.0$ & $-1.0$ & $1.00$\\
			4 & $3.8$ & $ 1.2$ & $1.44$\\
			5 & $4.6$ & $-0.6$ & $0.36$\\
			\midrule
			&       & $\sum e_i=0$ & $\text{SSE}=3.60$\\
			\bottomrule
		\end{tabular}
	\end{center}
	The residuals sum to zero---a general property of the least-squares fit
	(equation \eqref{eq:normal1}). The coefficient of determination is
	\[
	R^2=1-\frac{\text{SSE}}{S_{yy}}=1-\frac{3.6}{10}=0.64=r^2 ,
	\]
	confirming $R^2=r^2$ for simple linear regression.
	
	\begin{figure}[ht]
		\centering
		\begin{tikzpicture}
			\begin{axis}[
				width=11cm,height=7cm,
				xlabel={$x$ (size)},ylabel={$y$ (price)},
				xmin=0,xmax=6,ymin=0,ymax=6,
				grid=both,legend pos=north west]
				\addplot[only marks,mark=*,mark size=2.5pt,color=blue]
				coordinates {(1,1)(2,3)(3,2)(4,5)(5,4)};
				\addlegendentry{data}
				\addplot[thick,red,domain=0:6] {0.6+0.8*x};
				\addlegendentry{$\hat y=0.6+0.8x$}
				% residual segments
				\addplot[gray,dashed] coordinates {(1,1)(1,1.4)};
				\addplot[gray,dashed] coordinates {(2,3)(2,2.2)};
				\addplot[gray,dashed] coordinates {(3,2)(3,3.0)};
				\addplot[gray,dashed] coordinates {(4,5)(4,3.8)};
				\addplot[gray,dashed] coordinates {(5,4)(5,4.6)};
			\end{axis}
		\end{tikzpicture}
		\caption{Least-squares fit for the house-price dataset. Dashed segments are the
			residuals whose squared lengths are minimized.}
		\label{fig:fit}
	\end{figure}
	
	%==============================================================================
	\chapter{Complete Worked Example: Gradient Descent}
	\label{ch:gdexample}
	%==============================================================================
	
	We now fit the \emph{same} dataset with gradient descent, so we can compare the
	iterative answer against the exact closed-form solution
	$(\theta_0,\theta_1)=(0.6,0.8)$.
	
	\section{Setup}
	Model $\htheta(x)=\theta_0+\theta_1x$, cost
	$J=\frac{1}{2m}\sum(\htheta(x_i)-y_i)^2$ with $m=5$. Initialize
	$\theta_0=\theta_1=0$ and use learning rate $\alpha=0.1$. The update rules are
	\eqref{eq:gdupdate} with gradients
	\[
	g_0=\frac1m\sum_i(\htheta(x_i)-y_i),\qquad
	g_1=\frac1m\sum_i(\htheta(x_i)-y_i)\,x_i .
	\]
	
	\section{The Complete Iteration Table}
	Table~\ref{tab:gdfull} carries out the algorithm in full: for \emph{every}
	iteration it lists all five data points with their prediction
	$\hat y_i=\theta_0+\theta_1 x_i$, error $e_i=\hat y_i-y_i$, squared error
	$e_i^2$, and the product $e_i x_i$ used by the slope gradient. The shaded
	\emph{summary line} closing each iteration accumulates the column sums and
	turns them into the cost and the two gradients,
	\[
	J=\frac{1}{2m}\sum_i e_i^2=\frac{\sum e_i^2}{10},\qquad
	g_0=\frac{1}{m}\sum_i e_i=\frac{\sum e_i}{5},\qquad
	g_1=\frac{1}{m}\sum_i e_i x_i=\frac{\sum e_i x_i}{5},
	\]
	and finally applies the update
	$\theta_j:=\theta_j-\alpha\,g_j$ (with $\alpha=0.1$) to produce the
	$(\theta_0,\theta_1)$ that opens the next iteration. Reading the table top to
	bottom therefore reproduces the entire optimisation by hand.
	
	\begin{center}
		\footnotesize
		\setlength{\tabcolsep}{5pt}
		\renewcommand{\arraystretch}{1.15}
		\begin{longtable}{@{}c c c c c c r r r@{}}
			\caption{One large end-to-end trace of batch gradient descent on the
				house-price data ($\alpha=0.1$, five iterations). Each block shows the
				per-point computation; the summary line gives the sums, cost, gradients and the
				resulting parameter update.}\label{tab:gdfull}\\
			\toprule
			It & $\theta_0$ & $\theta_1$ & $x_i$ & $y_i$ & $\hat y_i$ & $e_i$ & $e_i^2$ & $e_i x_i$\\
			 &  &  &  &  & $\theta_0+\theta_1 x_i$ & $\hat y_i-y_i$ &  & \\
			\midrule
			\endfirsthead
			\multicolumn{9}{@{}l}{\footnotesize\itshape (Table~\ref{tab:gdfull} continued)}\\
			\toprule
			It & $\theta_0$ & $\theta_1$ & $x_i$ & $y_i$ & $\hat y_i$ & $e_i$ & $e_i^2$ & $e_i x_i$\\
			\midrule
			\endhead
			\midrule
			\multicolumn{9}{r@{}}{\footnotesize\itshape continued on next page}\\
			\endfoot
			\bottomrule
			\endlastfoot
			% ---------------- Iteration 0 ----------------
			\multirow{5}{*}{0} & \multirow{5}{*}{$0.0000$} & \multirow{5}{*}{$0.0000$}
			& 1 & 1 & $0.0000$ & $-1.0000$ & $1.0000$  & $-1.0000$\\
			& & & 2 & 3 & $0.0000$ & $-3.0000$ & $9.0000$  & $-6.0000$\\
			& & & 3 & 2 & $0.0000$ & $-2.0000$ & $4.0000$  & $-6.0000$\\
			& & & 4 & 5 & $0.0000$ & $-5.0000$ & $25.0000$ & $-20.0000$\\
			& & & 5 & 4 & $0.0000$ & $-4.0000$ & $16.0000$ & $-20.0000$\\
			\cmidrule(l){4-9}
			\multicolumn{6}{@{}r}{$\textstyle\sum$\;=} & $-15.0000$ & $55.0000$ & $-53.0000$\\
			 &  &  &  &  &  &  &  & \\
			\multicolumn{9}{@{}l}{\;\;$J=\tfrac{55.0000}{10}=5.5000,\quad
				g_0=\tfrac{-15.0000}{5}=-3.0000,\quad g_1=\tfrac{-53.0000}{5}=-10.6000
				\;\Rightarrow\; \theta_0{:}{=}0.3000,\ \theta_1{:}{=}1.0600$}\\
			\midrule
			It & $\theta_0$ & $\theta_1$ & $x_i$ & $y_i$ & $\hat y_i$ & $e_i$ & $e_i^2$ & $e_i x_i$\\
			 &  &  &  &  & $\theta_0+\theta_1 x_i$ & $\hat y_i-y_i$ &  & \\
			\midrule
			% ---------------- Iteration 1 ----------------
			\multirow{5}{*}{1} & \multirow{5}{*}{$0.3000$} & \multirow{5}{*}{$1.0600$}
			& 1 & 1 & $1.3600$ & $0.3600$  & $0.1296$ & $0.3600$\\
			& & & 2 & 3 & $2.4200$ & $-0.5800$ & $0.3364$ & $-1.1600$\\
			& & & 3 & 2 & $3.4800$ & $1.4800$  & $2.1904$ & $4.4400$\\
			& & & 4 & 5 & $4.5400$ & $-0.4600$ & $0.2116$ & $-1.8400$\\
			& & & 5 & 4 & $5.6000$ & $1.6000$  & $2.5600$ & $8.0000$\\
			\cmidrule(l){4-9}
			\multicolumn{6}{@{}r}{$\textstyle\sum$\;=} & $2.4000$ & $5.4280$ & $9.8000$\\
			 &  &  &  &  &  &  &  & \\
			\multicolumn{9}{@{}l}{\;\;$J=\tfrac{5.4280}{10}=0.5428,\quad
				g_0=\tfrac{2.4000}{5}=0.4800,\quad g_1=\tfrac{9.8000}{5}=1.9600
				\;\Rightarrow\; \theta_0{:}{=}0.2520,\ \theta_1{:}{=}0.8640$}\\
			\midrule
			It & $\theta_0$ & $\theta_1$ & $x_i$ & $y_i$ & $\hat y_i$ & $e_i$ & $e_i^2$ & $e_i x_i$\\
			 &  &  &  &  & $\theta_0+\theta_1 x_i$ & $\hat y_i-y_i$ &  & \\
			\midrule
			% ---------------- Iteration 2 ----------------
			\multirow{5}{*}{2} & \multirow{5}{*}{$0.2520$} & \multirow{5}{*}{$0.8640$}
			& 1 & 1 & $1.1160$ & $0.1160$  & $0.0135$ & $0.1160$\\
			& & & 2 & 3 & $1.9800$ & $-1.0200$ & $1.0404$ & $-2.0400$\\
			& & & 3 & 2 & $2.8440$ & $0.8440$  & $0.7123$ & $2.5320$\\
			& & & 4 & 5 & $3.7080$ & $-1.2920$ & $1.6693$ & $-5.1680$\\
			& & & 5 & 4 & $4.5720$ & $0.5720$  & $0.3272$ & $2.8600$\\
			\cmidrule(l){4-9}
			\multicolumn{6}{@{}r}{$\textstyle\sum$\;=} & $-0.7800$ & $3.7626$ & $-1.7000$\\
			 &  &  &  &  &  &  &  & \\
			\multicolumn{9}{@{}l}{\;\;$J=\tfrac{3.7626}{10}=0.3763,\quad
				g_0=\tfrac{-0.7800}{5}=-0.1560,\quad g_1=\tfrac{-1.7000}{5}=-0.3400
				\;\Rightarrow\; \theta_0{:}{=}0.2676,\ \theta_1{:}{=}0.8980$}\\
			\midrule
			It & $\theta_0$ & $\theta_1$ & $x_i$ & $y_i$ & $\hat y_i$ & $e_i$ & $e_i^2$ & $e_i x_i$\\
			 &  &  &  &  & $\theta_0+\theta_1 x_i$ & $\hat y_i-y_i$ &  & \\
			\midrule
			% ---------------- Iteration 3 ----------------
			\multirow{5}{*}{3} & \multirow{5}{*}{$0.2676$} & \multirow{5}{*}{$0.8980$}
			& 1 & 1 & $1.1656$ & $0.1656$  & $0.0274$ & $0.1656$\\
			& & & 2 & 3 & $2.0636$ & $-0.9364$ & $0.8768$ & $-1.8728$\\
			& & & 3 & 2 & $2.9616$ & $0.9616$  & $0.9247$ & $2.8848$\\
			& & & 4 & 5 & $3.8596$ & $-1.1404$ & $1.3005$ & $-4.5616$\\
			& & & 5 & 4 & $4.7576$ & $0.7576$  & $0.5740$ & $3.7880$\\
			\cmidrule(l){4-9}
			\multicolumn{6}{@{}r}{$\textstyle\sum$\;=} & $-0.1920$ & $3.7034$ & $0.4040$\\
			 &  &  &  &  &  &  &  & \\
			\multicolumn{9}{@{}l}{\;\;$J=\tfrac{3.7034}{10}=0.3703,\quad
				g_0=\tfrac{-0.1920}{5}=-0.0384,\quad g_1=\tfrac{0.4040}{5}=0.0808
				\;\Rightarrow\; \theta_0{:}{=}0.2714,\ \theta_1{:}{=}0.8899$}\\
			\midrule
			It & $\theta_0$ & $\theta_1$ & $x_i$ & $y_i$ & $\hat y_i$ & $e_i$ & $e_i^2$ & $e_i x_i$\\
			 &  &  &  &  & $\theta_0+\theta_1 x_i$ & $\hat y_i-y_i$ &  & \\
			\midrule
			% ---------------- Iteration 4 ----------------
			\multirow{5}{*}{4} & \multirow{5}{*}{$0.2714$} & \multirow{5}{*}{$0.8899$}
			& 1 & 1 & $1.1614$ & $0.1614$  & $0.0260$ & $0.1614$\\
			& & & 2 & 3 & $2.0513$ & $-0.9487$ & $0.9001$ & $-1.8974$\\
			& & & 3 & 2 & $2.9412$ & $0.9412$  & $0.8859$ & $2.8236$\\
			& & & 4 & 5 & $3.8311$ & $-1.1689$ & $1.3663$ & $-4.6755$\\
			& & & 5 & 4 & $4.7210$ & $0.7210$  & $0.5199$ & $3.6052$\\
			\cmidrule(l){4-9}
			\multicolumn{6}{@{}r}{$\textstyle\sum$\;=} & $-0.2940$ & $3.6981$ & $0.0172$\\
			 &  &  &  &  &  &  &  & \\
			\multicolumn{9}{@{}l}{\;\;$J=\tfrac{3.6981}{10}=0.3698,\quad
				g_0=\tfrac{-0.2940}{5}=-0.0588,\quad g_1=\tfrac{0.0172}{5}=0.0034
				\;\Rightarrow\; \theta_0{:}{=}0.2773,\ \theta_1{:}{=}0.8896$}\\
		\end{longtable}
	\end{center}
	
	\noindent\textbf{How to read one iteration (take iteration~0).}
	All parameters start at zero, so every prediction $\hat y_i=0$ and each error
	is simply $e_i=-y_i$. Squaring and summing the errors gives
	$\sum e_i^2=55$, hence $J=55/10=5.5$. The two gradients are the column means
	$g_0=\sum e_i/5=-3$ and $g_1=\sum e_i x_i/5=-10.6$. Stepping opposite the
	gradient, $\theta_0=0-0.1(-3)=0.3$ and $\theta_1=0-0.1(-10.6)=1.06$, which are
	exactly the values heading iteration~1. Every subsequent block repeats this
	identical recipe.
	
	\section{Discussion}
	Within a single step the cost collapses from $5.5$ to $0.54$, and the slope
	$\theta_1$ locks onto its optimal value $\approx0.89$ within a few iterations.
	The intercept $\theta_0$, however, drifts only slowly toward its true value
	$0.6$. This is the classic symptom of \emph{unscaled features}: because $x$
	ranges over $1$--$5$, the cost bowl is highly elongated and gradient descent
	zig-zags. Standardizing $x$ first would let both parameters converge together
	in a handful of steps. Given enough iterations the trajectory converges to the
	closed-form optimum $(0.6,\,0.8)$ with minimum cost $J=3.6/10=0.36$.
	
	\begin{figure}[ht]
		\centering
		\begin{tikzpicture}
			\begin{axis}[
				width=11cm,height=6.5cm,
				xlabel={Iteration},ylabel={Cost $J(\bm\theta)$},
				xmin=0,xmax=5,ymin=0,ymax=6,
				grid=both,mark size=2pt]
				\addplot[thick,blue,mark=*]
				coordinates {(0,5.5)(1,0.5428)(2,0.37626)(3,0.37034)(4,0.36981)(5,0.36948)};
			\end{axis}
		\end{tikzpicture}
		\caption{Cost versus iteration for gradient descent. The rapid initial drop and
			subsequent flattening are typical of convergence toward the global minimum.}
		\label{fig:costiter}
	\end{figure}
	
	%==============================================================================
	\chapter{MATLAB Implementation}
	%==============================================================================
	
	The following script reproduces both solutions and the two figures.
	
	\begin{lstlisting}[style=matlab,caption={Linear regression in MATLAB.}]
% ----- Data -----
x = [1;2;3;4;5];      % feature (column vector)
y = [1;3;2;5;4];      % target
m = length(y);
X = [ones(m,1), x];   % design matrix with intercept column

% ----- 1) Closed-form (normal equation) -----
theta_cf = (X' * X) \ (X' * y);   % '\' solves without explicit inverse
fprintf('Closed form:  theta0 = %.4f, theta1 = %.4f\n', ...
theta_cf(1), theta_cf(2));

% ----- 2) Gradient descent -----
theta = [0; 0];
alpha = 0.1;
iters = 50;
Jhist = zeros(iters,1);
for k = 1:iters
	h    = X * theta;                 % predictions
	err  = h - y;                     % residuals
	grad = (1/m) * (X' * err);        % gradient vector
	theta = theta - alpha * grad;     % simultaneous update
	Jhist(k) = (1/(2*m)) * (err' * err);
end
fprintf('Grad descent: theta0 = %.4f, theta1 = %.4f\n', ...
theta(1), theta(2));

% ----- 3) Cost vs iteration -----
figure; plot(1:iters, Jhist, 'b-o','LineWidth',1.2);
xlabel('Iteration'); ylabel('J(\theta)'); title('Cost vs Iteration'); grid on;

% ----- 4) Regression line -----
figure; scatter(x, y, 60, 'filled'); hold on;
xx = linspace(0, 6, 100);
plot(xx, theta_cf(1) + theta_cf(2)*xx, 'r-', 'LineWidth', 1.5);
xlabel('x'); ylabel('y'); legend('data','fit'); title('Regression Line'); grid on;
	\end{lstlisting}
	
	%==============================================================================
	\chapter{Python Implementation}
	%==============================================================================
	
	\section{Using NumPy and Matplotlib}
	\begin{lstlisting}[style=python,caption={Closed form and gradient descent with NumPy.}]
import numpy as np
import matplotlib.pyplot as plt

# ----- Data -----
x = np.array([1, 2, 3, 4, 5], dtype=float)
y = np.array([1, 3, 2, 5, 4], dtype=float)
m = len(y)
X = np.column_stack([np.ones(m), x])       # design matrix

# ----- 1) Closed form (normal equation) -----
theta_cf = np.linalg.solve(X.T @ X, X.T @ y)   # stable, avoids explicit inverse
print(f"Closed form:  theta0={theta_cf[0]:.4f}, theta1={theta_cf[1]:.4f}")

# ----- 2) Gradient descent -----
theta = np.zeros(2)
alpha, iters = 0.1, 50
J_hist = []
for _ in range(iters):
		err  = X @ theta - y
		grad = (X.T @ err) / m
		theta -= alpha * grad
		J_hist.append((err @ err) / (2 * m))
		print(f"Grad descent: theta0={theta[0]:.4f}, theta1={theta[1]:.4f}")

# ----- 3) Plots -----
plt.figure(); plt.plot(J_hist, 'b-o'); plt.xlabel("iteration")
plt.ylabel("J"); plt.title("Cost vs Iteration"); plt.grid(True)

plt.figure(); plt.scatter(x, y, label="data")
xx = np.linspace(0, 6, 100)
plt.plot(xx, theta_cf[0] + theta_cf[1]*xx, 'r', label="fit")
plt.xlabel("x"); plt.ylabel("y"); plt.legend(); plt.grid(True)
plt.show()
	\end{lstlisting}
	
	\section{Using Scikit-learn}
	\begin{lstlisting}[style=python,caption={The same fit in three lines with scikit-learn.}]
from sklearn.linear_model import LinearRegression
import numpy as np

x = np.array([1, 2, 3, 4, 5]).reshape(-1, 1)   # column of features
y = np.array([1, 3, 2, 5, 4])

model = LinearRegression().fit(x, y)
print("intercept theta0 =", model.intercept_)   # ~0.6
print("slope     theta1 =", model.coef_[0])      # ~0.8
print("R^2 =", model.score(x, y))                # 0.64
	\end{lstlisting}
	
	Both implementations return $\theta_0\approx0.6$, $\theta_1\approx0.8$, matching
	the hand computation of Chapter~\ref{ch:cfexample}.
	
	%==============================================================================
	\chapter{Comparison: Closed Form vs Gradient Descent}
	%==============================================================================
	
	\begin{center}
		\renewcommand{\arraystretch}{1.25}
		\begin{tabular}{@{}p{5.4cm}p{5.4cm}@{}}
			\toprule
			\textbf{Closed Form (Normal Equation)} & \textbf{Gradient Descent}\\
			\midrule
			Exact solution in one shot.            & Iterative approximation.\\
			No learning rate to tune.              & Must choose learning rate $\alpha$.\\
			Requires the matrix inverse
			$(\Xmat^{\!\top}\Xmat)^{-1}$.           & No matrix inversion needed.\\
			Cost $\mathcal{O}(n^3)$ in the number
			of features $n$.                       & Cost $\mathcal{O}(k\,mn)$ for $k$
			iterations; scales well.\\
			Fast and ideal for small feature sets. & Suitable for very large datasets and
			many features.\\
			Fails/unstable if $\Xmat^{\!\top}\Xmat$
			is singular.                           & Still works (with care) via variants
			and regularization.\\
			\bottomrule
		\end{tabular}
	\end{center}
	
	\paragraph{Rule of thumb.} Use the normal equation when the number of features
	is small (say $n\lesssim 10^4$); switch to (stochastic) gradient descent for
	high-dimensional or streaming data.
	
	%==============================================================================
	\chapter{Exercises}
	%==============================================================================
	
	\section{Solved Problems}
	
	\begin{example}[Manual fit]
		Fit a line to $x=(0,1,2,3)$, $y=(1,2,3,4)$.\\
		\textbf{Solution.} $\xbar=1.5,\ \ybar=2.5$. Deviations of $x$: $-1.5,-0.5,0.5,1.5$;
		of $y$: $-1.5,-0.5,0.5,1.5$. Thus $S_{xy}=2.25+0.25+0.25+2.25=5$ and
		$S_{xx}=5$, so $\theta_1=5/5=1$ and $\theta_0=2.5-1(1.5)=1$. Fit:
		$\hat y=1+x$ (a perfect fit, $R^2=1$).
	\end{example}
	
	\begin{example}[Effect of an outlier]
		Take the dataset of Chapter~\ref{ch:cfexample} but change $y_4$ from $5$ to
		$50$. Qualitatively, what happens to the slope? \\
		\textbf{Solution.} The point $(4,50)$ has a large positive deviation
		$y_4-\ybar$, injecting a large positive term into $S_{xy}$ and pulling the
		slope sharply upward. Squared-error loss is highly sensitive to outliers; a
		robust loss (e.g.\ Huber) would resist this.
	\end{example}
	
	\begin{example}[One gradient-descent step]
		For $x=(1,2)$, $y=(1,2)$, $\theta=(0,0)$, $\alpha=0.1$, compute one update.\\
		\textbf{Solution.} $m=2$. Errors $e=(-1,-2)$, $g_0=(-1-2)/2=-1.5$,
		$g_1=(-1\cdot1-2\cdot2)/2=-2.5$. Update:
		$\theta_0=0-0.1(-1.5)=0.15$, $\theta_1=0-0.1(-2.5)=0.25$.
	\end{example}
	
	\begin{example}[Proof: line passes through the centroid]
		Show the least-squares line satisfies $\ybar=\theta_0+\theta_1\xbar$.\\
		\textbf{Solution.} This is exactly the rearranged first normal equation
		\eqref{eq:theta0}, obtained by setting $\partial S/\partial\theta_0=0$.
	\end{example}
	
	\begin{example}[$R^2=r^2$]
		Verify $R^2=r^2$ for the Chapter~\ref{ch:cfexample} data.\\
		\textbf{Solution.} $r=0.8\Rightarrow r^2=0.64$; and
		$R^2=1-\text{SSE}/S_{yy}=1-3.6/10=0.64$. \checkmark
	\end{example}
	
	\section{Unsolved Problems}
	
	\begin{exercise}[Manual calculation]
		Fit $\hat y=\theta_0+\theta_1x$ to $x=(2,4,6,8)$, $y=(3,7,5,10)$ using the
		covariance/variance formulas. Report $\theta_0,\theta_1$ and $R^2$.
	\end{exercise}
	
	\begin{exercise}[Prediction]
		Using the model $\hat y=0.6+0.8x$ from Chapter~\ref{ch:cfexample}, predict the
		price for a house of size $x=6$ and state one reason such extrapolation may be
		unreliable.
	\end{exercise}
	
	\begin{exercise}[Learning-rate experiment]
		For the dataset of Chapter~\ref{ch:gdexample}, run gradient descent with
		$\alpha\in\{0.01,\,0.1,\,0.4\}$. For which value does the cost diverge? Explain.
	\end{exercise}
	
	\begin{exercise}[Derivation]
		Starting from $J(\vw)=\frac{1}{2m}(\Xmat\vw-\vy)^{\!\top}(\Xmat\vw-\vy)$, derive
		$\nabla J(\vw)=\frac1m\Xmat^{\!\top}(\Xmat\vw-\vy)$ and hence the normal
		equation.
	\end{exercise}
	
	\begin{exercise}[Proof-based]
		Prove that $\Xmat^{\!\top}\Xmat$ is positive semi-definite and give a condition
		under which it is positive definite (hence invertible).
	\end{exercise}
	
	\begin{exercise}[Conceptual]
		Explain why the mean squared error is preferred over the mean absolute error
		when a smooth, differentiable optimization landscape is desired, and name one
		situation where mean absolute error is nonetheless preferable.
	\end{exercise}
	
	\begin{exercise}[Feature scaling]
		Standardize the feature $x$ of Chapter~\ref{ch:gdexample} to zero mean and unit
		variance, re-run five iterations of gradient descent, and compare the
		convergence speed of $\theta_0$ and $\theta_1$ with the unscaled case.
	\end{exercise}
	
	\begin{exercise}[Coding]
		Implement ridge regression by modifying the normal equation to
		$\vw=(\Xmat^{\!\top}\Xmat+\lambda\mathbf{I})^{-1}\Xmat^{\!\top}\vy$ (do not
		penalize the intercept) and study the effect of $\lambda$ on the fitted slope.
	\end{exercise}
	
	%==============================================================================
	\appendix
	\chapter{Notation Summary}
	%==============================================================================
	\begin{center}
		\begin{tabular}{@{}ll@{}}
			\toprule
			Symbol & Meaning\\
			\midrule
			$m$ & number of training examples\\
			$n$ & number of features\\
			$x^{(i)},\,y^{(i)}$ & feature(s) and target of the $i$-th example\\
			$\vw=(\theta_0,\dots,\theta_n)$ & parameter (weight) vector\\
			$\htheta(\vx)$ & hypothesis / prediction\\
			$\Xmat$ & design matrix, $\R^{m\times(n+1)}$\\
			$\vy$ & target vector, $\R^{m}$\\
			$\mathbf{r}=\Xmat\vw-\vy$ & residual vector\\
			$J(\vw)$ & mean-squared-error cost\\
			$\alpha$ & learning rate\\
			$S_{xx},S_{xy}$ & $\sum(x_i-\xbar)^2$, $\sum(x_i-\xbar)(y_i-\ybar)$\\
			$r,\ R^2$ & correlation coefficient, coefficient of determination\\
			\bottomrule
		\end{tabular}
	\end{center}
	
\end{document}
